\begin{center}
	\zihao{4}\heiti
	%
	\LARGE{ 中国科学院大学\\ $\boldsymbol{2021}$年招收攻读硕士学位研究生入学统一考试试卷\\
	科目名称：电动力学}

\end{center}
\noindent
考生须知：\\
1.本试卷满分为150分,全部考试时间总计180分钟。\\
2.所有答案必须写在答题纸上,写在试题纸上或草稿纸上一律无效。

\vspace*{2ex}
\hrule
\vspace*{4ex}

\fancypagestyle{mypagestyle}{%
	\fancyhf{}% Clear header/footer
	\fancyfoot[OR]{第~\thepage~页，共2页}%
	\fancyfoot[EL]{第~\thepage~页，共2页}%
	\fancyfoot[OL]{科目名称：电动力学 }%
	\fancyfoot[ER]{科目名称：电动力学 }%
	\renewcommand{\headrulewidth}{0pt}% Header rule of .4pt
	\renewcommand{\footrulewidth}{0.5pt}% Header rule of .4pt
}
\fancypagestyle{plain}{%
	\fancyhf{}% Clear header/footer
	\fancyfoot[OR]{第~\thepage~页，共2页}%
	\fancyfoot[EL]{第~\thepage~页，共2页}%
	\fancyfoot[OL]{科目名称：电动力学 }%
	\fancyfoot[ER]{科目名称：电动力学 }%
}

\setcounter{page}{1}%设置页面计数器
\pagestyle{mypagestyle}

\onehalfspacing
% 排版结束，下面是题目

\vspace{0.4cm}

\noindent
一、简答题（共40分）
\vspace{0.2cm}

\noindent
1.\hspace{0.8em}从麦克斯韦方程出发，推导洛伦兹规范下的达朗贝尔方程
\vspace{0.1cm}

\noindent
2.\hspace{0.8em}写出电偶极距的表达式，并用它表示出电势
\vspace{0.1cm}

\noindent
3.\hspace{0.8em}写出磁偶极距的表达式，并用它表示出磁标势和磁矢量
\vspace{0.4cm}

\noindent
二、\hspace{0.2em}单项选择题（共30分，每小题5分）
\vspace{0.2cm}

\noindent
1.\hspace{0.8em}$\varepsilon$与$\mu$的单位是
\vspace{0.2cm}

A.\hspace{0.8em}[安]$^{-2}$ [千克]$^1 $[秒]$^{-4}$ [米]$^{3}$，[安]$^2$ [千克]$^{-1}$ [秒]$^2$ [米]$^{-1}$\vspace{0.15cm}

B.\hspace{0.8em}[库伦]$^{-2}$[千克]$^{1}$[秒]$^{0}$[米]$^{1}$， [库伦]$^{2}$[千克]$^{-1}$[秒]$^{2}$[米]$^{-3}$ \vspace{0.15cm}

C.\hspace{0.8em}[安]$^{2}$ [千克]$^{-1} $[秒]$^{4}$ [米]$^{-3}$，[安]$^{-2}$ [千克]$^{1}$ [秒]$^{-2}$ [米]$^{1}$\vspace{0.15cm}

D.\hspace{0.8em}[库伦]$^{-2}$[千克]$^{1}$[秒]$^{-2}$[米]$^{3}$， [库伦]$^{2}$[千克]$^{-1}$[秒]$^{0}$[米]$^{-1}$ \vspace{0.15cm}

E.\hspace{0.8em}[安]$^2$ [千克]$^{-1}$ [秒]$^2$ [米]$^{-1}$，[安]$^{-2}$ [千克]$^1 $[秒]$^{-4}$ [米]$^{3}$

\vspace{0.2cm}

\noindent
2.\hspace{0.8em}磁场在分界面的边界条件是
\vspace{0.2cm}

A.\hspace{0.8em}$\Vec{e_n} \times (\Vec{B_2} - \Vec{B_1}) = \alpha_f$
\hspace{5em}
B.\hspace{0.8em}$\Vec{e_n} \times (\Vec{H_2} - \Vec{H_1}) = \alpha_f$\vspace{0.2cm}

C.\hspace{0.8em}$\nabla \times (\Vec{B_2} - \Vec{B_1}) = \alpha_f$
\hspace{5.1em}
D.\hspace{0.8em}$\nabla \times (\Vec{H_2} - \Vec{H_1}) = \alpha_f$
\vspace{0.2cm}

\noindent
3.\hspace{0.8em}电磁波入射到良导体内部，以下说法不正确的是\vspace{0.2cm}

A.\hspace{0.8em}透射强度几乎为零\hspace{11em}
B.\hspace{0.8em}穿透深度几乎为零\vspace{0.1cm}

C.\hspace{0.8em}界面处总电场强度几乎为零\hspace{7em}
D.\hspace{0.8em}界面处的电流密度几乎为零
\vspace{0.2cm}

\noindent
4.\hspace{0.8em}对于电荷均匀分布的椭球，以下正确的是
\vspace{0.2cm}

A.\hspace{0.8em}电偶极距和电四极距都为零\hspace{8em}
B.\hspace{0.8em}电偶极距为零，电四极距不为零\vspace{0.1cm}

C.\hspace{0.8em}电偶极距不为零，电四极距为零\hspace{6em}
D.\hspace{0.8em}电偶极矩和电四极矩都不为零\vspace{0.2cm}

\noindent
5.\hspace{0.8em}一场源线度为$l$，电磁波的波长为$\lambda$，场源在远处距离为$R$处场点产生辐射时，对所需要做的近似和一些表述正确的是
\vspace{0.2cm}

A.\hspace{0.8em}只需要近似到$\displaystyle{\frac{1}{R}}$，且$l \ll R$，$R \ll \lambda$，辐射强度正比于$\displaystyle{(\frac{l}{\lambda})^2}$\vspace{0.2cm}

B.\hspace{0.8em}只需要近似到$\displaystyle{\frac{1}{R}}$，且$l \gg R$，$R \gg \lambda$，辐射强度正比于$\displaystyle{(\frac{l}{\lambda})^2}$\vspace{0.2cm}

C.\hspace{0.8em}只需要近似到$\displaystyle{\frac{1}{R^2}}$，且$l \ll R$，$R \ll \lambda$，辐射强度正比于$\displaystyle{(\frac{l}{\lambda})^4}$\vspace{0.2cm}

D.\hspace{0.8em}只需要近似到$\displaystyle{\frac{1}{R^2}}$，且$l \gg R$，$R \gg \lambda$，辐射强度正比于$\displaystyle{(\frac{l}{\lambda})^4}$

\newpage

\noindent
6.\hspace{0.8em}一个带电的粒子已接近于光速的速度运动，则它产生的电场
\vspace{0.2cm}

A.\hspace{0.8em}集中在与运动方向垂直的方向\hspace{5em}
B.\hspace{0.8em}集中在与运动方向平行的方向\vspace{0.1cm}

C.\hspace{0.8em}均匀分布\hspace{14em}
D.\hspace{0.8em}逐渐减小趋于零\vspace{1cm}

\noindent
三、\hspace{0.2em}（共20分）电磁波从介质1入射到介质2，两介质都有$\mu = \mu_0$，电场垂直于入射面\vspace{0.2cm}

\noindent
1、写出电场和磁场在分界面的边界条件并证明:\vspace{0.2cm}

$\displaystyle{ \frac{ E^{'} }{ E } = - \frac{ \sin (\theta - \theta^{''} ) }{ \sin (\theta + \theta^{''} ) } } \quad \displaystyle{ \frac{ E^{''} }{ E } =  \frac{ 2 \cos \theta \sin \theta^{''} ) }{ \sin (\theta + \theta^{''} ) } }$\vspace{0.2cm}

\noindent
2、假设介质1为真空，介质2中$\varepsilon_r = 2$，入射角$\theta = 45^{\circ}$，求反射系数和折射系数。\vspace{3cm}

\noindent
四、\hspace{0.2em}（共20分）一通电流$I$的通电螺线管，单位长度所绕匝数为n，管半径为R\vspace{0.2cm}

\noindent
1、求管内外的磁感应强度$\Vec{B}$\vspace{0.2cm}

\noindent
2、求管内外的磁矢量$\Vec{A}$
\vspace{3cm}

\noindent
五、\hspace{0.2em}（共20分）一电荷量为q的带电粒子垂直于磁场以速$V_0$入射，磁场强度为$B_0$
\vspace{0.2cm}

\noindent
1、计算该带电粒子辐射的电偶极势和电磁场
\vspace{0.2cm}

\noindent
2、计算瞬间能流密度和瞬间功率\vspace{3cm}


\noindent
六、\hspace{0.2em}（共20分）参考系$k^{'}$相对参考系$k$运动速度为$v$。
\vspace{0.2cm}

\noindent
1、写出两参考系中的时空变换关系
\vspace{0.2cm}

\noindent
2、计算出三维速度矢量关系
\vspace{0.2cm}

\noindent
3、写出四维速度$U_{\mu}$的定义，并计算它的洛伦兹不变式
